    \documentclass[11pt]{article}

    \usepackage{fullpage}
    \usepackage{graphicx}
    \usepackage{times}
    \usepackage{url}
    \usepackage{booktabs}

    \title{\bf Smart Course Registration and Scheduling System}
    \author{
        \bf Lohit P Talavar \qquad 210564 \qquad lohitpt21@iitk.ac.in
    }
    \date{CS315 Final Project --- Indian Institute of Technology Kanpur}

    \begin{document}

    \maketitle

    \begin{abstract}
    In this project I have designed and implemented a relational database system
    for managing university course registrations using PostgreSQL. The schema consist
    of six main entities which are Students, Courses, Instructors, Departments,
    Enrollments and Prerequisites. Three real world business rules are enforced
    directly at the database layer itself — prerequisite checking, seat limit and
    schedule conflict detection. All tables are normalized upto Third Normal Form.
    Six indexes including a partial index on active enrollments are added to improve
    query performance. The project covers all major topics of CS315 like ER modeling,
    normalization, SQL, constraints, transactions, triggers, stored procedures and views.
    \end{abstract}

    % ------------------------------------------------------------------
    \section{Motivation and Problem Statement}
    % ------------------------------------------------------------------

    Course registration in a university involves lot of rules and constraints that
    needs to be handled properly. A student cannot enroll in a course without first
    completing its prerequisite courses. Also the course should not be full and there
    should be no clash in the timetable with other courses the student is already
    enrolled in.

    Most of the time these rules are checked in the application code which means
    any client that directly access the database can bypass them easily. The main
    goal of this project was to push all these three constraints inside the database
    itself using triggers and check constraints so that no one can bypass them.

    Apart from correctness the system also needs to support efficient queries for
    common operations like timetable look ups, roster retrieval, enrollment validation
    and occupancy reporting. This domain directly maps to the CS315 syllabus topics ---
    ER modeling, normalization, indexing, transactions and views --- all in one single
    coherent system. No external APIs or complex real time data was needed, everything
    can be done purely in PostgreSQL which made it a good choice.

    % ------------------------------------------------------------------
    \section{Methodology}
    % ------------------------------------------------------------------

    I followed a milestone based approach where I first did the design and then
    started implementing each part one by one.

    \paragraph{ER Design.}
    I identified six entities and mapped out all the relationships before writing
    any SQL. Some important design decisions were made at this stage. First I decided
    to use a separate \texttt{Prerequisites} table instead of a single \texttt{prereq\_id}
    column in the Courses table, this is because a single column can only store one
    prerequisite but a course like CS401 requires both CS315 and MTH201 so that would
    not work. Second I stored \texttt{enrolled\_count} as a column in Courses table
    itself (which is technically denormalization) so that the seat limit CHECK constraint
    can fire instantly, and a trigger keeps this count updated. Third the timeslot was
    stored as three separate columns (day, start time, end time) instead of a single
    string so that SQL can do proper time comparison for conflict detection.

    \paragraph{Normalization.}
    All tables were checked and proven to be in Third Normal Form by analyzing functional
    dependencies. The most important violations I avoided were storing \texttt{dept\_name}
    in the Students table (this would be a transitive dependency) and storing instructor
    details in the Courses table. Foreign keys are used instead everywhere.

    \paragraph{Schema and Sample Data.}
    The DDL defines all six tables with primary keys, foreign keys, NOT NULL and UNIQUE
    constraints. For the seed data I added 3 departments, 4 instructors, 8 courses,
    10 students and 22 enrollments. I also made sure to include edge cases like a full
    course, a course with prerequisites, and a pair of courses with same timeslot.

    \paragraph{Business Logic and Constraints.}
    A single BEFORE INSERT trigger handles all the three business rule checks before
    any row is inserted into the Enrollments table. It reads the Courses row for seat
    check, queries Enrollments for completed prerequisites, and checks time overlap
    with the students existing active courses.

    \paragraph{Indexing.}
    Six indexes were added based on which columns appear most in WHERE clauses.
    I used EXPLAIN ANALYZE to compare query plans before and after adding indexes.

    \paragraph{Transactions and Triggers.}
    A separate AFTER trigger maintains the \texttt{enrolled\_count} after every insert,
    update or delete on Enrollments. The enrollment transaction uses SELECT FOR UPDATE
    to lock the course row so that two students cannot grab the same last seat
    simultaneously. I also wrote a stored procedure \texttt{enroll\_student()} that
    wraps the whole flow.

    \paragraph{Views and Reporting.}
    Three views are created for common reporting needs. Six complex analytical queries
    were also written covering occupancy rates, credit loads and instructor workload.

    % ------------------------------------------------------------------
    \section{Implementation and Results}
    % ------------------------------------------------------------------

    The complete code is available at \url{https://github.com/lohitpt252003/CS315}.

    \paragraph{Schema.}
    The \texttt{schema.sql} file defines all six tables. Some of the important
    constraints are shown below:

    \begin{verbatim}
    CONSTRAINT chk_capacity_limit CHECK (enrolled_count <= capacity)
    CHECK (status IN ('enrolled','dropped','waitlisted','completed'))
    CONSTRAINT no_self_prereq     CHECK (course_id <> prereq_course_id)
    UNIQUE (student_id, course_id)
    \end{verbatim}

    \paragraph{Business Rule Testing.}
    Table~\ref{tab:tests} shows the test cases I ran against the seed data to verify
    all three constraints are working correctly.

    \begin{table}[h]
    \centering
    \begin{tabular}{lll}
    \toprule
    \textbf{Test} & \textbf{Expected} & \textbf{Outcome} \\
    \midrule
    4th student in CS201 (capacity 3) & Rejected & course full (3/3) \\
    Charlie enrolls in CS315 (no prereq) & Rejected & prerequisite not met \\
    Charlie enrolls in EE201 (timeslot clash) & Rejected & schedule conflict \\
    Charlie enrolls in EE301 (all ok) & Accepted & row inserted successfully \\
    Frank drops EE301 & enrolled\_count $-1$ & trigger verified \\
    \bottomrule
    \end{tabular}
    \caption{Constraint test results using seed data.}
    \label{tab:tests}
    \end{table}

    \paragraph{Indexing Results.}
    Among all six indexes, the partial index on active enrollments was the most useful
    one:

    \begin{verbatim}
    CREATE INDEX idx_enrollments_enrolled_partial
        ON Enrollments(student_id, course_id)
        WHERE status = 'enrolled';
    \end{verbatim}

    This index is used on every enrollment attempt by both the prerequisite check and
    the conflict check. It is also smaller than a full index since it excludes
    completed and dropped rows. Before adding indexes the planner was doing sequential
    scans on Enrollments. After indexes it uses Index Scan which is significantly faster
    when the table grows large.

    \paragraph{Transaction and ACID.}
    The enrollment transaction shows all four ACID properties. Atomicity means the
    insert and count update happens together or not at all. Consistency is maintained
    by the trigger. Isolation is handled by FOR UPDATE lock which prevents two
    transactions from reading same available seat. Durability is ensured by COMMIT.

    \paragraph{Views Output.}
    Some sample outputs from the three views with seed data:
    \begin{itemize}
    \item \texttt{student\_transcript} for Alice: shows CS201 (MWF 10:00--11:00) and
            MTH101 (MWF 11:00--12:00).
    \item \texttt{course\_roster} for CS201: 3 students enrolled, 0 seats remaining.
    \item \texttt{department\_summary}: CS department has 4 courses with total 8 enrollments.
    \end{itemize}

    % ------------------------------------------------------------------
    \section{Discussion and Limitations}
    % ------------------------------------------------------------------

    There are few limitations in the current implementation that should be noted.

    \paragraph{Day pattern matching.}
    The schedule conflict detection is using exact string comparison on
    \texttt{schedule\_day}. So two courses like \texttt{'MWF'} and \texttt{'MW'} would
    not be detected as conflicting even thought they share Monday and Wednesday. A proper
    solution would be to use bitmask representation for days. This is a known limitation
    and was kept simple intentionally.

    \paragraph{Concurrency bottle neck.}
    Every enrollment locks the Courses row to update enrolled\_count. If many students
    try to register for same popular course at same time, they will all be queued one
    after another. At large scale this can become a serious performance bottleneck.
    Some solutions like optimistic locking or queue based enrollment system would help
    here.

    \paragraph{No archiving.}
    The Enrollments table keeps growing with no way to archive old data. Completed
    enrollments from past semesters slow down queries that filter by active status.
    Partitioning by semester would be the right solution for this.

    \paragraph{Timeslot model simplification.}
    Currently each course is assumed to meet on same set of days every week. Courses
    with irregular schedules like alternating weeks lab sessions cannot be stored in
    this schema without changing it.

    % ------------------------------------------------------------------
    \section{Contributions}
    % ------------------------------------------------------------------

    Lohit P Talavar did everything including the ER design, schema definition,
    normalization proof, seed data, all SQL files for constraints, indexing, triggers,
    transactions, stored procedures, views and analytical queries, and also wrote this
    report.

    \end{document}
